\documentclass[11pt]{article}

\usepackage[margin=1in]{geometry}
\usepackage{amsmath,amssymb}
\usepackage{booktabs}
\usepackage{natbib}
\usepackage{url}

\newcommand{\magmaan}{\texttt{magmaan}}
\newcommand{\lavaan}{\texttt{lavaan}}
\newcommand{\Cov}{\operatorname{Cov}}
\newcommand{\Var}{\operatorname{Var}}

\title{Fast start-value portfolios for hard small-sample structural equation models}
\author{ }
\date{\today}

\begin{document}
\maketitle

\begin{abstract}
Iterative estimation of structural equation models can fail because the sample
is small, the model is weakly measured, the parameter space has nearby
boundaries, or the optimizer begins in an unhelpful region. Recent bounded
random-start proposals show that even ordinary SEMs can benefit from multiple
starting points. This note studies a complementary question: under a fixed
small compute budget, can a portfolio of deterministic, bounded-random, and
cheaply screened starting values improve convergence, admissibility, and basin
selection for hard linear SEMs? We use \magmaan{} to generate and fit large
grids of small-sample CFA and SEM designs, recording optimizer diagnostics,
positive-definiteness failures, objective values, gradients, wall time, and
agreement with heavier multistart references.
\end{abstract}

\section{Motivation}

Default SEM starting values remain modest. For measurement models, software may
combine simple constants with noniterative estimators such as FABIN
\citep{haggund1982fabin}, Guttman's multiple-group method
\citep{guttman1952multiple}, or Bentler's noniterative CFA estimator
\citep{bentler1982noniterative}. For structural regressions and covariances,
zeros remain common. Recent work by \citet{dejonckere2025randomstarts}
demonstrates that bounded random starting values can rescue many small-sample
SEM fits that default starts miss.

This paper asks whether a start-value portfolio can do better than any one
scheme alone. The target is not a universal closed-form start. The target is a
practical policy that spends a small, explicit budget on candidate generation,
screening, and fitting.

\section{Contribution}

The planned contribution has four parts.

\begin{enumerate}
\item Reproduce published hard small-sample designs from
  \citet{dejonckere2022bounded}, \citet{dejonckere2025randomstarts}, and
  \citet{ludtke2021pmlmcmc}.
\item Build a corpus of genuinely hard cases for modern optimizers, including
  failures, boundary solutions, non-PD trajectories, large final gradients, and
  disagreements across start policies.
\item Compare deterministic starts, bounded random starts, and small start
  portfolios under equal budgets.
\item Separate convergence from solution quality by reporting admissibility,
  objective value, stationarity, standard-error availability, and agreement
  with heavy multistart references.
\end{enumerate}

\section{Start Policies}

The initial policy set will include simple starts, FABIN2/FABIN3, Guttman,
Bentler-1982, James-Stein-style CFA starts, bounded random starts, and
hybrids. A candidate can be scored before full optimization by the discrepancy
value, implied covariance positive-definiteness margin, projected bounds
distance, and a short safeguarded optimization probe.

\section{Simulation Designs}

\subsection{Small latent regression}

The first design follows the two-factor latent regression used by
\citet{dejonckere2022bounded} and \citet{dejonckere2025randomstarts}. Each
latent is measured by three indicators, with weak enough measurement and small
enough samples to create frequent nonconvergence and improper solutions.

\subsection{Crossloading chain SEM}

The second design is the three-factor chain SEM, adapted from
\citet{chen2001improper}, used in both De Jonckere and Rosseel papers. It has
three indicators per factor, three crossloadings, and a structural chain
$\eta_1 \rightarrow \eta_2 \rightarrow \eta_3$. This design is useful because
the number and placement of omitted crossloadings can be varied while holding
the data-generating model fixed.

\subsection{Two-factor weak-loading CFA}

The third design follows \citet{ludtke2021pmlmcmc}: two correlated factors,
three indicators per factor, standardized loadings around $.50$, small sample
sizes, and latent correlations ranging from moderate to high. This design is a
good admissibility and boundary stress test.

\subsection{Multistate-single-trait model}

The fourth design follows the multistate-single-trait example from
\citet{dejonckere2025randomstarts}. The first implementation uses the reported
state/trait loading and measurement-error values, with explicit nuisance
defaults for the remaining population quantities. This section will be tightened
against the authors' OSF scripts before results are reported.

\section{Outcomes}

For each fit we will record convergence status, optimizer status, function and
gradient evaluations, final gradient norm, objective value, positive-definiteness
failures, variance/correlation admissibility, standard-error availability, wall
time, and agreement with a heavier reference run.

\section{Open Questions}

\begin{itemize}
\item Which failures remain hard for \magmaan{} with NLopt L-BFGS?
\item Are start portfolios still useful when the optimizer almost always
  returns a solution?
\item Does a start policy improve admissibility, or merely increase the number
  of converged-but-improper fits?
\item Can cheap screening predict the best start well enough to avoid full
  multistart fitting?
\end{itemize}

\bibliographystyle{apalike}
\bibliography{convergence-note}

\end{document}
